%% ============================================================
%%  Objetivos
%%  Duenio inicial: claude-1
%%  Reglas: ver ../../CLAUDE.md
%% ============================================================

\chapter{Objetivos}
\label{cap:objetivos}

\section{Objetivo general}

Desarrollar un framework metodológico que, operando de forma agnóstica sobre cualquier modelo cuya atención se calcule entre tokens de variables con nombre, transforme esa representación en una hipótesis relacional compacta e interpretable sobre las variables del fenómeno, y establecer el criterio que decide cuándo esa hipótesis puede sostenerse.

\section{Objetivos específicos}

Los tres objetivos siguientes cubren, en orden, la formalización del método, la construcción del criterio con que se lo evalúa y su aplicación sobre el caso de estudio.

\begin{enumerate}
\item Formalizar matemáticamente las transformaciones que llevan de la matriz de interacción entre variables al grafo dirigido, y de este a la hipótesis relacional compacta, junto con las condiciones que debe cumplir una representación para ser entrada válida del framework.

\item Diseñar un criterio que evalúe las relaciones extraídas y determine cuáles se sostienen como hipótesis estructurales genuinas.

\item Aplicar el framework sobre el caso de estudio de índices espectrales y evaluar la hipótesis obtenida con el criterio anterior, reportando relación por relación cuáles se sostienen y cuáles quedan descartadas.
\end{enumerate}
